\documentclass[11pt,a4paper]{article}
\usepackage[T1]{fontenc}
\usepackage{amsmath,amssymb,amsthm,mathtools}
\usepackage{geometry}
\usepackage[hidelinks]{hyperref}
\geometry{margin=1in}

\newtheorem{theorem}{Theorem}[section]
\newtheorem{lemma}[theorem]{Lemma}
\newtheorem{proposition}[theorem]{Proposition}
\newtheorem{corollary}[theorem]{Corollary}
\theoremstyle{definition}
\newtheorem{definition}[theorem]{Definition}
\newtheorem{example}[theorem]{Example}

\author{Part II}
\date{Recorded from the lecturer in 2026-2027.}

\title{Asymptotic Methods}
\begin{document}
\maketitle
\thispagestyle{empty}

\noindent\textbf{Pre-requisites}
\begin{itemize}
    \item Complex methods / analysis.
    \item Refer to Section II in Stuart.
\end{itemize}

\section*{1. Intro \& Motivation}

\begin{itemize}
    \item Most real problems don't have closed form solutions, or if they do, they are unworkable or difficult to interpret.

    \item Numerical approximations can be useful, but they can be expensive, they may have bad numerical properties, or solutions may be inscrutable.

    \item Taylor series, etc.\ can be limited in applicability.

    \item Asymptotic methods emphasise usefulness, interpretability as well as accuracy.
\end{itemize}

\noindent\textbf{Examples}

\medskip

\noindent\textbf{e.g. 1)}
\[
\Gamma(z)=\int_0^\infty t^{z-1}e^{-t}\,dt,
\qquad \operatorname{Re}(z)>0.
\]

For $n$ an integer,
\[
\Gamma(n+1)=n!
\sim \sqrt{2\pi}\,n^{n+\frac12}e^{-n}
\left(
1+\frac{1}{12n}+\frac{1}{288n^2}+\cdots
\right)
\qquad \text{as }n\to\infty.
\]

\medskip

\noindent\textbf{e.g. 2)}
\[
x^2y''+xy'-x^2y=0.
\]

Exact series solution:
\[
I_0(x)=\sum_{n=0}^{\infty}
\frac{x^{2n}}{2^{2n}(n!)^2}.
\]

Liouville--Green:
\[
\Longrightarrow
I_0(x)\sim(2\pi x)^{-1/2}e^x
\qquad \text{as }x\to\infty.
\]

\medskip

\noindent\textbf{e.g. 3)}
\[
y''+k^2e^{-x^2/2}y=0,
\]
with
\[
y\to y_0 \qquad \text{as }x\to\infty,
\]
\[
y\to \pm y_0 \qquad \text{as }x\to-\infty.
\]

\[
k=?,
\qquad k=k_0,k_1,k_2,\ldots
\]

\[
\Longrightarrow
k_n\sim\frac{\sqrt{\pi}}{2}
\left(n+\frac12\right)
\qquad \text{as }n\to\infty.
\]

\noindent WKBJ.

\medskip

\begin{itemize}
    \item Convergent methods typically give arbitrarily accurate approximations with more \& more terms/iterations.

    \item Asymptotic methods give arbitrarily accurate approximations as some variable approaches a limit. Thus they always come with a qualifier, e.g.\ $z\to0$ or $n\to\infty$.
\end{itemize}


\section*{2. Asymptotic Expansions}

\subsection*{2.1 Definitions \& Notation}

\subsubsection*{(a) Asymptotic Equality/Equivalence}

For 2 functions $f(x)$ and $g(x)$, we have
\[
f(x)\sim g(x)
\qquad \text{about }x=x_0
\]
if
\[
\boxed{
\lim_{x\to x_0}\frac{f(x)}{g(x)}=1.
}
\]

\noindent\textbf{Remarks:}

\begin{itemize}
    \item Usually we want to approximate $f$ with $g$.

    \item We can write
    \[
    f(x)=
    \underbrace{g(x)}_{\text{dominant}}
    +
    \underbrace{[f(x)-g(x)]}_{\text{subdominant}}.
    \]

    \item Note $\sim$ is an equivalence relation
    (that is reflexive, transitive, symmetric).

    \item It respects products:
    \[
    f_1\sim g_1,\qquad f_2\sim g_2
    \]
    \[
    \Longrightarrow f_1f_2\sim g_1g_2.
    \]

    \item It is linear:
    \[
    \alpha f_1+\beta f_2
    \sim \alpha g_1+\beta g_2,
    \]
    where $\alpha,\beta$ are constants.

    \item Definition has problems when dealing with oscillatory functions that go through $0$ as $x\to x_0$.
\end{itemize}


\subsubsection*{(b) ``Little $o$''}

Suppose $f$ \& $g$ are real functions.

We write
\[
\boxed{
f(x)=o(g(x))
\qquad \text{as }x\to x_0
}
\]
if
\[
\boxed{
\lim_{x\to x_0}\frac{f(x)}{g(x)}=0.
}
\]

\noindent\textbf{Remark:}

\[
f(x)=o(g(x))
\]
can also be written as
\[
f(x)\ll g(x)
\qquad \text{as }x\to x_0.
\]

\begin{itemize}
    \item Can be framed in inequalities, e.g.\ $x_0=\infty$ and let $g>0$.

    Then $f=o(g)$ as $x\to\infty$ is equivalent to
    \[
    \forall\varepsilon>0,\quad
    \exists a
    \quad\text{s.t. when }x>a,
    \]
    \[
    \boxed{|f|\leq\varepsilon g.}
    \]
\end{itemize}


\subsubsection*{(c) ``Big $O$''}

We write
\[
f(x)=O(g(x))
\qquad \text{as }x\to x_0
\]
for $g>0$, if
\[
\exists C>0
\quad\text{s.t.}
\]
\[
\boxed{
|f(x)|\leq Cg(x)
}
\]
for $x$ sufficiently close to $x_0$.

\medskip

\noindent\textbf{Warning:}

In some areas of physics people use big-$O$ and $\sim$ interchangeably.


\subsubsection*{Little Exercises: Show That}

\noindent\textbf{(a)}
\[
f(x)\sim g(x)
\iff
f-g=o(g)
\qquad \text{as }x\to x_0.
\]

\medskip

\noindent\textbf{(b)}
\[
\sin x\ll x^2
\qquad \text{as }x\to\infty.
\]

\medskip

\noindent\textbf{(c)}
\[
12x^5+6x^3+x=O(x^5)
\qquad \text{as }x\to\infty.
\]

\medskip

\noindent\textbf{(d)}
\[
e^{ax}\ll e^{bx}
\qquad \text{as }x\to\infty,
\]
if $b>a>0$.

\medskip

\noindent\textbf{(e)}
\[
e^x\gg x^N
\qquad \text{as }x\to\infty,
\]
\[
\forall N\in\mathbb{Z}^+.
\]


\subsection*{2.2 Asymptotic Expansions}

\begin{itemize}
    \item Have unworkable/unknown $f(x)$.

    \item Want to know its behaviour in a limit.

    \item Want to find a simpler/known $g(x)$ so that
    \[
    f\sim g.
    \]

    \item Also want to break up $g$ into a hierarchy of approximations, s.t.
    \[
    g(x)=\sum_{k=0}^{N}a_k\phi_k(x),
    \]
    where $a_k$ are coefficients, $\phi_k$ are known \& workable, and where adding more terms gives a better approximation that can be estimated.
\end{itemize}


\subsubsection*{(a) Asymptotic Sequence}

A sequence of functions
\[
\{\phi_0(x),\phi_1(x),\ldots\}
\]
is an ``asymptotic sequence'' as $x\to x_0$ if
\[
\boxed{
\phi_{k+1}(x)\ll\phi_k(x)
\qquad \text{as }x\to x_0,
\quad \forall k.
}
\]

\medskip

\noindent\textbf{e.g. 1)}
\[
\phi_k=(x-a)^k
\qquad \text{around }x=a.
\]

Check:
\[
\frac{\phi_{k+1}}{\phi_k}
=x-a\to0
\qquad \text{as }x\to a.
\quad \checkmark
\]

\medskip

\noindent\textbf{e.g. 2)}
\[
\phi_k=e^{-kx}x^k
\qquad \text{as }x\to\infty
\text{ or }x\to-\infty.
\]

Check:
\[
\frac{\phi_{k+1}}{\phi_k}
=e^{-x}x.
\]

\[
\begin{aligned}
e^{-x}x&\to0
&&\text{as }x\to\infty
&&\checkmark\\
e^{-x}x&\to\infty
&&\text{as }x\to-\infty
&&\times
\end{aligned}
\]

\medskip

\noindent\textbf{e.g. 3)}
\[
\phi_k=\cos(nx)
\qquad \text{as }x\to\infty.
\]

\[
\times
\]

\medskip

\noindent\textbf{e.g. 4)}
\[
\left\{
|x|,\,
\log|x|,\,
\log|\log|x||,\,
\ldots
\right\}
\qquad \text{as }x\to\infty.
\]

\[
?
\]


\subsubsection*{(b) Asymptotic Expansion}

\begin{itemize}
    \item Suppose $f(x)$ is a function we want to understand near $x=x_0$.

    \item Let $\{\phi_k\}_{k=0}^{\infty}$ be an asymptotic sequence.
\end{itemize}

Then, the finite sum
\[
\sum_{k=0}^{N}a_k\phi_k(x)
\]
for some coefficients, is an ($N$-term) asymptotic approximation of $f(x)$ as $x\to x_0$ if:
\[
\boxed{
f(x)=\sum_{k=0}^{N}a_k\phi_k(x)+o(\phi_N(x))
}
\qquad \text{as }x\to x_0.
\]

Alternatively, we can write the above as
\[
\boxed{
\frac{f(x)-\displaystyle\sum_{k=0}^{N}a_k\phi_k(x)}
{\phi_N(x)}
\longrightarrow 0
}
\qquad \text{as }x\to x_0.
\]

\begin{center}
    \textbf{End of Lecture 1}
\end{center}
\end{document}
